\documentclass[11pt,a4paper]{article}

\usepackage[utf8]{inputenc}
\usepackage[T1]{fontenc}
\usepackage[british]{babel}
\usepackage[margin=2.5cm]{geometry}
\usepackage{parskip}
\usepackage{enumitem}
\usepackage{titlesec}
\usepackage[hidelinks]{hyperref}

% Section spacing tuned for a short policy-style document
\titlespacing*{\section}{0pt}{1.4em}{0.6em}
\titlespacing*{\subsection}{0pt}{1.0em}{0.4em}

\setlist[itemize]{itemsep=0.2em, topsep=0.3em}

\title{ACTS Contributor Policy \\ \large Contributorship, Authorship, and Citation}
\author{}
\date{\docstamp}

% Document status and revision stamp; see actsdoc.sty.
\usepackage[status=draft, version=0.1, name={ACTS Contributor Policy}]{actsdoc}

\begin{document}

\maketitle
\thispagestyle{empty}

\section*{Purpose}

This document defines who counts as an ACTS contributor and how that determination
is made. The contributor list serves two purposes: it is the public record of who
has shaped the project, and it seeds the author list for ``on behalf of ACTS''
publications and talks. The policy is intentionally light, leaving room for the
collaborative judgement that has served the project well, while making the
underlying mechanics legible to people inside and outside the project.

The contributor list is maintained as a two-tier list: \emph{active} and
\emph{inactive}. Active contributors seed the author list for new project releases,
the pool from which CODEOWNERS are drawn, and the seed for ``on behalf of ACTS''
author lists. Both active and inactive contributors are listed on the project's
contributor page.

\section*{Sufficient conditions for active contributor status}

A person qualifies as an active contributor if they meet \emph{any one} of the
following conditions within the past 12 months. The list is meant to be inclusive:
meeting one condition is enough; meeting several does not confer additional status.
The project assumes good faith and collaborative intent throughout --- these
conditions describe the kinds of contribution that earn active status, not a
checklist for adjudication.

\begin{itemize}
    \item \textbf{Code-based engagement.} Authored or substantively reviewed at
          least two substantial pull requests on the project's code repositories.
          Authorship and review are treated as fungible toward this threshold. A
          pull request is "substantial" if it makes non-trivial changes to the
          library itself, fundamental infrastructure or documentation. A review 
          is substantive if it goes beyond a bare approval ---
          comments, change requests, or threaded discussion all count. Self-reviews
          do not count. Reviews on pull requests that did not merge still count.
    \item \textbf{Project authorship.} Contributed substantively to an ACTS
          publication --- writing, reviewing the manuscript, producing results or
          figures, contributing the code or analysis the publication describes, or
          similar work on the publication itself. Appearing on an author list by
          virtue of contributor status, without doing work on the publication, does
          not satisfy this condition.
    \item \textbf{Supervision.} Supervised a student or postdoc whose work was
          substantively on ACTS, where the supervisee is themselves a contributor
          under one of the other conditions. The supervision has to relate to the 
          ACTS contribution itself and not to unrelated activities of the supervisee.
    \item \textbf{Sustained personnel contribution.} Made it possible for one or
          more individuals to spend significant amount of time substantially
          working on ACTS development, integration, or maintenance. The
          individuals doing the work qualify on their own merits via the
          conditions above; this condition recognizes the person who enabled the
          sustained engagement.
    \item \textbf{Sustained material contribution.} Provided computing
          infrastructure, testing resources, hosting, event funding, or similar
          in-kind contributions on which the project materially depended.
\end{itemize}

Giving a talk on behalf of ACTS is not by itself a sufficient condition for
contributor status --- see the section on representing the project below. Holding
the role of Project Contact or experiment liaison is likewise not a sufficient
condition: PCs in practice qualify under the code-based or supervision conditions
many times over, and liaisons represent their experiment to ACTS rather than
contributing to ACTS itself, so the role is orthogonal to contributor status. As
with talks, liaisons are often contributors via the other conditions; this is
correlation, not consequence.

\section*{Active and inactive}

Active contributors are those who have met a sufficient condition within the past
12 months. Inactive contributors are those who were once active but have not met a
condition in the past 12 months. Movement between tiers is automatic and is not a
judgement of the person; inactive status reflects current involvement, not the value
of past contributions. Inactive contributors who become active again return to the
active list at the next quarterly update.

The active list seeds the author list for new project releases and the project's
CODEOWNERS pool. CODEOWNERS membership is by self-selection: active contributors who
agree to be code owners for specific paths are listed; nobody is added to CODEOWNERS
without their consent.

Both active and inactive contributors are listed on the project's contributor page,
with their status clearly indicated. Past project releases preserve the author list
as it was at the time of release, so contributions made during a person's active
period remain permanently associated with them through the historical record.

\section*{How contributors are recognised}

The contributor list is updated by pull request, opened against the project's
source-of-truth file. The Project Contact opens a regular update PR on a roughly
quarterly cadence:

\begin{itemize}
    \item The PR is \textbf{seeded mechanically} from project history, drawing on
          commit and review activity in the lookback window.
    \item Conditions that are not visible in the project's code history --- project
          authorship, supervision, sustained personnel contribution, sustained
          material contribution --- are reflected through either
          \textbf{self-nomination} or manual addition by the PC. Anyone who
          believes they meet such a condition can be added in two equivalent
          ways:
          \begin{itemize}
              \item Asking the PC to include them in the next quarterly PR, with a
                    brief description of the contribution. Nominations on someone
                    else's behalf are welcome, including from the experiment liaison
                    group.
              \item Opening a PR themselves against the source-of-truth file, with
                    the brief description of the contribution in the PR body.
          \end{itemize}
    \item The PR is \textbf{announced at the weekly developer meeting} when
          opened, and circulated on the project's mailing list and main Mattermost
          channel.
    \item The PR remains open for \textbf{one week} for objections, comments, and
          additions. Silence is consent.
    \item After the objection window closes, the PR is merged.
\end{itemize}

Three of the four quarterly updates handle additions: contributors who newly meet
a sufficient condition are added. The fourth quarterly update, scheduled to align
with the annual workshop, additionally handles the active-to-inactive transition
for contributors who have not met a sufficient condition in the past 12 months.

Contributors can be added at any time during a quarter; the quarterly cadence is
the default rhythm, not a constraint. Self-nomination PRs and other interim updates
are merged outside the quarterly schedule when ready. If a publication or talk is
being prepared and the active list needs updating ahead of the next quarterly PR,
the PC can also open an interim PR.

The project assumes good faith. In practice, mechanical seeding plus self-nomination
is expected to produce a list that the community recognises as accurate without
further intervention. If a concern is raised --- typically through the objection
window on a PR --- the PC facilitates a brief discussion to address it. Concerns
are expected to be rare; the document does not prescribe a formal adjudication
process because none has been needed.

\section*{Representing the project}

The project does not maintain a process for approving who may present on behalf of
ACTS, and contributor status is not a hard requirement for doing so. In practice
this question rarely creates difficulty, and the project prefers not to formalise
what has not been a problem.

The expectation is that anyone intending to give a talk, write a proceedings, or
otherwise represent ACTS publicly \textbf{circulates the intent} in advance ---
typically at the developer meeting, via the mailing list or Mattermost or
directly to the PC. There is no explicit approval step. If concerns are raised
about the planned representation, they can be voiced to the PC, who facilitates
a consensus discussion in the usual way. In the great majority of cases,
circulation is a courtesy that lets the project know what is happening and gives
others the opportunity to coordinate or contribute material.

Giving a talk on behalf of ACTS does not by itself confer contributor status, just
as contributor status is not required to give such a talk. The two are deliberately
decoupled.

\section*{Authorship for ``on behalf of ACTS'' publications and talks}

The active contributor list, at the time a publication or talk is being prepared,
seeds the author list for work presented on behalf of ACTS. The project does not
formalise what makes a publication or talk ``on behalf of ACTS'' --- that
determination is made collaboratively in the same way the project handles other
questions.

People who become active during preparation can be added through the normal
mechanism; people who transition to inactive during preparation remain on the
author list because they were active when the work happened. Individuals listed as
active may decline authorship on any specific publication.

Active contributor status is the seed for the author list, not a strict membership
rule. People who are not on the active list may be appropriate co-authors for a
specific publication, and people on the active list may decline. The active list
provides a starting point that captures the community at the time of writing
without being treated as exclusive.

\section*{Citing ACTS}

The recommended citation depends on what your collaboration's policies allow:

\begin{itemize}
    \item \textbf{Preferred:} The 2021 paper [Ai et al., \texttt{arXiv:2106.13593}]
          together with the Zenodo DOI for the specific version of ACTS used in the
          work. This combines the canonical peer-reviewed reference with a
          reproducible, historically accurate snapshot of who was active when the
          cited version was released. The version-specific DOI is available from
          the project's Zenodo record at
          \href{https://doi.org/10.5281/zenodo.5141418}{\texttt{doi.org/10.5281/zenodo.5141418}}
          --- follow the link to the relevant release.
    \item \textbf{Alternative}, where peer-reviewed citation is not required: the
          Zenodo DOI for the specific version used, together with the Zenodo
          concept DOI (\texttt{10.5281/zenodo.5141418}).
    \item \textbf{If your collaboration's policies do not permit citing software
          DOIs:} the 2021 paper alone.
\end{itemize}

Citing only the concept DOI is discouraged: the concept DOI resolves to the latest
release, so the author list it points to changes over time, which obscures who
actually contributed to the version used.

A suggested acknowledgement phrasing, where a full citation is not appropriate (for
example in theses):

\begin{quote}
\itshape
This work made use of A Common Tracking Software (ACTS) version vXX.Y.Z [Ai et al.,
\texttt{2106.13593}; Zenodo: \texttt{10.5281/zenodo.NNNNNNNN}]. We thank the ACTS
developers for their continued support of the toolkit.
\end{quote}

The project will also maintain a public list of papers that have used ACTS,
populated by self-report from authors. This gives contributors a visible record of
the project's impact.

\section*{Bookkeeping}

The project maintains its citation and contributor metadata in a single
source-of-truth file in the main code repository, from which downstream artefacts
are generated automatically. The relevant artefacts are:

\begin{itemize}
    \item The \textbf{Zenodo metadata} used for new releases, which lists active
          contributors as authors. Past releases preserve the author list as it was
          at the time of release.
    \item The \textbf{AUTHORS file} in the repository, which lists all contributors
          grouped by active and inactive status.
    \item The \textbf{CODEOWNERS configuration}, which is drawn from active
          contributors who have opted in to specific paths.
\end{itemize}

Updates to the source-of-truth file flow through the standard pull request
mechanism described above. Generated files are kept in sync automatically, so
contributors normally do not need to edit them directly. The exact technical
arrangement of source-of-truth file and synchronisation tooling is an
implementation matter and is documented separately in the repository.

\section*{Review}

This policy will be reviewed annually, alongside the PC and ELG terms of reference,
and adjusted as practice settles. The thresholds and exclusion list in particular
are expected to evolve based on experience.

\end{document}